\chapter{Architecture réseau et infrastructure}

\section{Vue d'ensemble}

L'infrastructure corpmsi.local est composée de 6 VLANs isolés, hébergés sur une machine physique exécutant VMware Workstation Pro. Chaque VLAN représente un département ou service.

\section{Diagramme réseau complet}

\begin{figure}[H]
    \centering
    \fbox{\parbox{0.85\linewidth}{\centering
    \begin{verbatim}
┌─────────────────────────────────────────────────────┐
│         Internet / WAN                              │
│       (réseau hôte / NAT VirtualBox)                │
└──────────────────┬──────────────────────────────────┘
                   │
                   │ em0 WAN (DHCP)
          ┌────────────────────┐
          │    pfSense CE      │
          │  192.168.40.1/24   │
          │   Firewall/Router  │
          └────┬───────────────┘
               │ em1–em6 (6 VLANs)
    ┌──────────┼──────────────────────┐
    │          │                      │
┌───────┐  ┌────────┐  ┌────────┐  ┌────────┐
│VLAN10 │  │VLAN20  │  │VLAN30  │  │VLAN40  │  ┌─────────┐ ┌──────────┐
│DIR    │  │COMM    │  │IT      │  │SERVERS │  │VLAN50   │ │VLAN60    │
│em6    │  │em2     │  │em3     │  │em1     │  │RH (em4) │ │COMPT(em5)│
└───────┘  └────────┘  └────────┘  └────────┘  └─────────┘ └──────────┘
    \end{verbatim}
    \vspace{0.5cm}}}
    \caption{Topologie réseau complète}
    \label{fig:topology}
\end{figure}

\section{Adressage IP — Plan complet}

\begin{longtable}{|p{3cm}|p{3cm}|p{3cm}|p{2cm}|}
    \hline
    \textbf{VLAN} & \textbf{Réseau} & \textbf{Passerelle} & \textbf{DHCP} \\
    \hline
    \endhead
    
    V10 Direction & 192.168.10.0/24 & 192.168.10.1 & Oui \\
    \hline
    V20 Commercial & 192.168.20.0/24 & 192.168.20.1 & Non \\
    \hline
    V30 Informatique & 192.168.30.0/24 & 192.168.30.1 & Oui \\
    \hline
    V40 Serveurs & 192.168.40.0/24 & 192.168.40.1 & Non \\
    \hline
    V50 RH & 192.168.50.0/24 & 192.168.50.1 & Oui \\
    \hline
    V60 Comptabilité & 192.168.60.0/24 & 192.168.60.1 & Oui \\
    \hline
    
    \caption{Plan d'adressage des VLANs}
    \label{tab:vlans}
\end{longtable}

\section{Adresses IP fixes}

\begin{table}[H]
    \centering
    \begin{tabularx}{\textwidth}{|l|l|p{4cm}|}
        \hline
        \textbf{Hôte} & \textbf{IP} & \textbf{Rôle} \\
        \hline
        srv-dc01 & 192.168.40.10 & DC principal, PDC, DNS, DHCP \\
        \hline
        srv-dc02 & 192.168.40.11 & DC secondaire, DNS, GC \\
        \hline
        srv-fichiers01 & 192.168.40.20 & Serveur fichiers (SMB, DFS) \\
        \hline
        pfSense (LAN) & 192.168.40.1 & Passerelle VLAN 40 \\
        \hline
        pfSense (em0) & DHCP / FAI & Accès Internet \\
        \hline
    \end{tabularx}
    \caption{Adresses IP fixes de l'infrastructure}
    \label{tab:ips}
\end{table}

\section{Structure Active Directory}

\begin{figure}[H]
    \centering
    \fbox{\parbox{0.85\linewidth}{\centering
    \begin{verbatim}
corpmsi.local (domaine)
├── OU=Direction
│   ├── alice.dupont
│   ├── GRP_Direction (groupe)
│   └── (5 postes clients)
├── OU=Commercial
│   ├── bob.martin
│   ├── charlie.bernard
│   ├── GRP_Commercial (groupe)
│   └── (3 postes clients)
├── OU=Informatique
│   ├── david.laurent
│   ├── GRP_Informatique (groupe)
│   └── (3 postes technicians)
├── OU=RH
│   ├── eve.moreau
│   ├── frank.petit
│   ├── GRP_RH (groupe)
│   └── (2 postes)
├── OU=Comptabilite
│   ├── grace.roux
│   ├── henry.blanc
│   ├── GRP_Comptabilite (groupe)
│   └── (2 postes)
└── OU=Serveurs
    ├── CN=srv-dc01
    ├── CN=srv-dc02
    └── CN=srv-fichiers01
    \end{verbatim}
    \vspace{0.5cm}}}
    \caption{Structure des Unités Organisationnelles}
    \label{fig:ou-structure}
\end{figure}

\section{Configuration VMware Workstation Pro}

\subsection{Virtual Network Editor}

\begin{table}[H]
    \centering
    \begin{tabularx}{\textwidth}{|l|l|l|l|}
        \hline
        \textbf{VMnet} & \textbf{Type} & \textbf{Sous-réseau} & \textbf{Utilisation} \\
        \hline
        VMnet0 & Bridged & (réseau physique) & WAN pfSense \\
        \hline
        VMnet1 & Host-only & 192.168.10.0/24 & V10 Direction \\
        \hline
        VMnet2 & Host-only & 192.168.20.0/24 & V20 Commercial \\
        \hline
        VMnet3 & Host-only & 192.168.30.0/24 & V30 Informatique \\
        \hline
        VMnet4 & Host-only & 192.168.40.0/24 & V40 Serveurs \\
        \hline
        VMnet5 & Host-only & 192.168.50.0/24 & V50 RH \\
        \hline
        VMnet6 & Host-only & 192.168.60.0/24 & V60 Comptabilité \\
        \hline
    \end{tabularx}
    \caption{Configuration VMware Workstation Pro}
    \label{tab:vmware}
\end{table}

\textbf{Important :} Le DHCP VMware doit être \textbf{désactivé} sur VMnet1–VMnet6. Le DHCP sera géré par les Domain Controllers via DHCP Relay depuis pfSense.

\subsection{Affectation réseau des VMs}

\begin{table}[H]
    \centering
    \begin{tabularx}{\textwidth}{|l|l|p{3cm}|}
        \hline
        \textbf{VM} & \textbf{Cartes réseau} & \textbf{Assignation} \\
        \hline
        pfSense & 8 cartes & em0=VMnet0 WAN, em1–em6=VMnet1–6 \\
        \hline
        srv-dc01 & 1 carte & VMnet4 (VLAN 40 Serveurs) \\
        \hline
        srv-dc02 & 1 carte & VMnet4 (VLAN 40 Serveurs) \\
        \hline
        srv-fichiers01 & 1 carte & VMnet4 (VLAN 40 Serveurs) \\
        \hline
        PC-BUR01 & 1 carte & VMnet1–6 (configurable pour test) \\
        \hline
    \end{tabularx}
    \caption{Affectation réseau des VMs}
    \label{tab:vm-network}
\end{table}

\section{Pile technique détaillée}

\begin{table}[H]
    \centering
    \begin{tabularx}{\textwidth}{|p{2.5cm}|p{1.5cm}|p{5.5cm}|}
        \hline
        \textbf{Outil} & \textbf{Version} & \textbf{Rôle dans l'infrastructure} \\
        \hline
        VMware Workstation Pro & 17.x & Hyperviseur pour toutes les VMs \\
        \hline
        Windows Server 2022 & Std & OS pour DCs et serveur fichiers \\
        \hline
        Windows 11 & 22H2 & Workstation client de test \\
        \hline
        pfSense CE & 2.7.x & Pare-feu, routeur, DHCP Relay \\
        \hline
        AD DS 2022 & — & Annuaire centralisé \\
        \hline
        DNS Server & 2022 & Résolution corpmsi.local \\
        \hline
        DHCP Server & 2022 & Attribution IP par VLAN \\
        \hline
        DFS Namespaces & 2022 & Partages centralisés transparents \\
        \hline
        Squid & 6.x & Proxy applicatif transparent \\
        \hline
        SquidGuard & 1.4.x & Filtrage URL + blacklists UT1 \\
        \hline
        OpenSSL & 3.x & SSL/TLS MITM pour inspection \\
        \hline
    \end{tabularx}
    \caption{Pile technique complète}
    \label{tab:stack}
\end{table}

\section{Département Direction — V10}

\begin{itemize}
    \item \textbf{Réseau} : 192.168.10.0/24
    \item \textbf{DHCP} : Oui (100–150)
    \item \textbf{Accès} : Serveurs AD + Internet (proxy filtrant)
    \item \textbf{Isolation} : Oui, séparé de V20/V50/V60
\end{itemize}

\section{Département Commercial — V20}

\begin{itemize}
    \item \textbf{Réseau} : 192.168.20.0/24
    \item \textbf{DHCP} : Non (adresses statiques)
    \item \textbf{Accès} : Serveurs AD + Internet (proxy filtrant)
    \item \textbf{Isolation} : Oui, adresses préassignées pour chaque vendeur
\end{itemize}

\section{Département Informatique — V30}

\begin{itemize}
    \item \textbf{Réseau} : 192.168.30.0/24
    \item \textbf{DHCP} : Oui (100–150)
    \item \textbf{Accès} : Tous les VLANs (administration complète)
    \item \textbf{Isolation} : Non — accès transversal complet requis
\end{itemize}

\section{Département Ressources Humaines — V50}

\begin{itemize}
    \item \textbf{Réseau} : 192.168.50.0/24
    \item \textbf{DHCP} : Oui (100–150)
    \item \textbf{Accès} : Serveurs AD + Internet restrictif (email uniquement)
    \item \textbf{Isolation} : Oui, données sensibles
\end{itemize}

\section{Département Comptabilité — V60}

\begin{itemize}
    \item \textbf{Réseau} : 192.168.60.0/24
    \item \textbf{DHCP} : Oui (100–150)
    \item \textbf{Accès} : Serveurs AD + Internet très restrictif (HTTPS only)
    \item \textbf{Isolation} : Oui, isolation maximale (données financières)
\end{itemize}

\section{Serveurs — V40}

\begin{itemize}
    \item \textbf{Réseau} : 192.168.40.0/24
    \item \textbf{DHCP} : Non (adresses statiques)
    \item \textbf{Accès} : Internet pour mises à jour, communication intra-DC
    \item \textbf{Isolation} : Bloqués d'initier vers les autres VLANs (sécurité)
\end{itemize}
